% Generated draft: requires mathematical and source review.
\documentclass[11pt]{article}
\usepackage[margin=25mm]{geometry}
\usepackage{fontspec}
\setmainfont{Latin Modern Roman}
\usepackage{amsmath,amssymb,mathtools,unicode-math}
\setmathfont{Latin Modern Math}
\usepackage{enumitem,needspace,xurl,hyperref}
\setlength{\parindent}{0pt}
\setlength{\parskip}{5pt}
\setlength{\emergencystretch}{3em}
\newcommand{\N}{\mathbb{N}}
\newcommand{\Z}{\mathbb{Z}}
\newcommand{\Q}{\mathbb{Q}}
\newcommand{\R}{\mathbb{R}}
\newcommand{\C}{\mathbb{C}}
\newcommand{\F}{\mathbb{F}}
\newcommand{\A}{\mathbb{A}}
\newcommand{\PP}{\mathbb{P}}
\newcommand{\E}{\mathbb{E}}
\newcommand{\eps}{\varepsilon}
\newcommand{\dd}{\,\mathrm{d}}
\begin{document}
\section*{UnsolvedMath ID: 701043}
\textbf{Source problem code:} catalogue entry
\subsection{Definitions and mathematical statement}
Let \(B\) be an abelian group, written additively, and let \(A\leq\operatorname{Aut}(B)\). A group is periodic if each of its elements has finite order. Interpret the source phrase ``group of regular automorphisms'' as requiring every nonidentity element of \(A\) to act without nonzero fixed points:
\[
\forall a\in A\setminus\{1\},\qquad \{b\in B:a(b)=b\}=\{0\}.
\]
The supplied record does not define ``regular''; this interpretation requires confirmation against the original source. No finiteness or torsion assumption on \(B\) is supplied.

\begin{enumerate}
\item[(a)] For every prime \(p\), if \(A\) is periodic, satisfies the displayed fixed-point condition, and has exponent exactly \(p\), must \(A\) be cyclic? Here exponent exactly \(p\) means that \(A\neq\{1\}\) and \(a^p=1\) for every \(a\in A\).
\item[(b)] If \(A\) is periodic, satisfies the displayed fixed-point condition, and admits a generating set \(S\subseteq A\) such that
\[
A=\langle S\rangle,\qquad \operatorname{ord}(s)=3\quad\text{for every }s\in S,
\]
must \(A\) be finite? The generating set is not assumed finite.
\end{enumerate}
The supplied record retains part (a) in the main problem section and marks part (b) as archived solved, reporting an affirmative answer attributed to A. Kh. Zhurtov. These are supplied status qualifications, not independently verified mathematical conclusions.
\subsection{Short English statement}
For a periodic group acting by fixed-point-free automorphisms on an abelian group, part (a) asks whether prime exponent forces cyclicity, and part (b) asks whether generation by elements of order \(3\) forces finiteness. The supplied record labels the overall problem partially solved and reports part (b) as affirmatively solved.
\subsection{Sources}
\begin{itemize}
\item Local UnsolvedMath registry, permanent ID \(701043\), exact source problem code \texttt{catalogue entry}, title \emph{Kourovka Notebook Problem 14.58}; supplied upstream revision \texttt{d371ee917ac9647fbb06cd6fc529deec6f3dda82}.
\item \emph{The Kourovka Notebook (Unsolved Problems in Group Theory)}, 14th issue (1999), Problem 14.58, proposed by V. D. Mazurov. The supplied record describes an update in the 21st issue (2026), identified as \texttt{arXiv:1401.0300v48}. The collection is attributed to editors E. I. Khukhro and V. D. Mazurov.
\item Supplied Archive of Solved Problems attribution for part (b): A. Kh. Zhurtov, \emph{Siberian Math. J.} \textbf{41}, no. 2 (2000), 268--275. The archive excerpt reports an affirmative answer; no article title is supplied.
\item The supplied status review is dated \texttt{2026-10-16}, labels the problem \texttt{partially\_solved}, gives its evidence as ``published'', and states that solved parts occur in the archive while remaining parts occur in the main section. This is supplied dated curation, not an independently verified literature assessment.
\end{itemize}
\end{document}
